\section{Metrics}

% Authors: Keerthi Rapolu & Sreeja Katta

PBCP is evaluated with a small metric set chosen to reflect operational goals
rather than metric diversity. CPS measures prevention value, ESR measures
whether prevention preserves useful execution, and IFS measures workload
alignment between intent and observed behavior.

\subsection{Cost Prevention Score}

\emph{Cost Prevention Score} (CPS) is the primary prevention metric. It measures
how much potential waste the system prevents before or during execution relative
to the relevant cost exposure. CPS therefore quantifies the benefit of
pre-billing intervention and runtime governance in operational rather than purely
predictive terms. In this paper, the denominator is always the \emph{unchanged
cost exposure}: for pre-provision experiments, the projected cost of executing
the submitted request as-is; for runtime experiments, the cost that the same
submitted configuration would have incurred without intervention, including any
avoidable idle or overrun interval. With that stage-specific denominator, CPS
remains bounded by 1.

\subsection{Execution Success Rate}

\emph{Execution Success Rate} (ESR) is the anti-gaming constraint. A governance
system can inflate prevention metrics by blocking useful work too aggressively,
so CPS is only meaningful when interventions still preserve successful workload
completion. ESR therefore distinguishes effective governance from trivial waste
avoidance through over-restriction. We therefore report \emph{Valid CPS} as
$\mathrm{CPS} \times \mathrm{ESR}$ so that high prevention is only credited when
execution remains operationally successful.

\subsection{Intent-Fit Score}

\emph{Intent-Fit Score} (IFS) is the diagnostic metric for workload alignment.
Within the intent-behavior divergence framing, the canonical IFS definition is
the cosine similarity between intent embeddings and behavior embeddings. In this
paper, any sub-score decomposition is interpretability support only, not the
primary metric definition. The relationship among the metrics is deliberate: CPS
indicates how much waste was prevented, ESR indicates whether prevention
remained operationally safe, and IFS helps explain why a workload diverged from
its expected behavior in the first place.
